% Generated draft: requires mathematical and source review.
\documentclass[11pt]{article}
\usepackage[margin=25mm]{geometry}
\usepackage{fontspec}
\setmainfont{Latin Modern Roman}
\usepackage{amsmath,amssymb,mathtools,unicode-math}
\setmathfont{Latin Modern Math}
\usepackage{enumitem,needspace,xurl,hyperref}
\setlength{\parindent}{0pt}
\setlength{\parskip}{5pt}
\setlength{\emergencystretch}{3em}
\newcommand{\N}{\mathbb{N}}
\newcommand{\Z}{\mathbb{Z}}
\newcommand{\Q}{\mathbb{Q}}
\newcommand{\R}{\mathbb{R}}
\newcommand{\C}{\mathbb{C}}
\newcommand{\F}{\mathbb{F}}
\newcommand{\A}{\mathbb{A}}
\newcommand{\PP}{\mathbb{P}}
\newcommand{\E}{\mathbb{E}}
\newcommand{\eps}{\varepsilon}
\newcommand{\dd}{\,\mathrm{d}}
\begin{document}
\section*{UnsolvedMath ID: 11000115}
\textbf{Source problem code:} catalogue entry
\subsection{Definitions and mathematical statement}
Let \(\phi(g,n)\) be the function defined immediately before Problem 2.10 in the source. Its definition, admissible integer ranges and codomain are absent from the supplied record. The requested tasks are:
\begin{enumerate}
\item Determine \(\phi(g,n)\) for its admissible parameters.
\item Decide whether it is constant, with the intended variable dependence made explicit.
\item If it is not constant, compare \(\phi(g,n)\) and \(\phi(h,n)\) whenever \(g<h\), holding \(n\) fixed.
\end{enumerate}
In particular, constancy in genus at fixed \(n\) would mean
\[\forall n\ \forall g,h\text{ admissible},\qquad \phi(g,n)=\phi(h,n),\]
whereas constancy on the full parameter domain also compares different \(n\). The source does not say which is intended. The comparison request specifies no inequality direction, so no increasing or decreasing monotonicity conjecture is inserted.

The research summary associates this with positive factorizations and Hurwitz-type counts, but gives no definition linking those topics to \(\phi\). It cannot justify selecting a count, a minimum, an equivalence relation on factorizations, or a finite/infinite-value convention. A precise numerical or enumerative statement requires the omitted definition; \(\phi\) here is an explicit unresolved source symbol, not a newly defined invariant.
\subsection{Short English statement}
Compute the source's function \(\phi(g,n)\), determine its constancy and compare its values as genus increases at fixed \(n\). The defining invariant is missing, so no counting interpretation or monotonicity is assumed.
\subsection{Sources}
\begin{itemize}
\item Local UnsolvedMath registry, permanent ID \(11000115\), exact source code \texttt{catalogue entry}.
\item Farb, \emph{Problems on Mapping Class Groups and Related Topics}, Problem 2.10, PDF page 96; the supplied running footer names M. Korkmaz: \url{https://www.math.uchicago.edu/~farb/papers/mcgbook.pdf}. Registry access 2026-07-29; pdf-font-structure extraction. No separate primary definition or result reference is supplied.
\end{itemize}
\end{document}
